\chapter{Datos}\label{cap:datos}

El trabajo usa dos fuentes. El \enquote{espacio} es el reposo en cama de la NASA: personas sanas acostadas durante
semanas con la cabeza inclinada a $-6^\circ$, medidas cada mañana por enfermería. La \enquote{Tierra} es la vida
diaria de personas con un reloj Fitbit, del conjunto público LifeSnaps, más los días previos al reposo de cada sujeto.

\section{El reposo en cama como análogo de la microgravedad}

En el reposo en cama con la cabeza hacia abajo, los fluidos se desplazan hacia el tórax y la cabeza y los músculos y
huesos dejan de cargar peso, como en el espacio. Por eso es el análogo terrestre más usado: \textcite{pavy2007}
revisaron veinte años de estudios de reposo en cama y lo que enseñaron de la fisiología humana, y
\textcite{hargens2016} discutieron por qué el reposo prolongado sirve como análogo de la microgravedad.

Los datos de este trabajo vienen del proyecto Flight Analogs de la NASA, en el Johnson Space Center y la University of
Texas Medical Branch (UTMB), publicados en el Life Sciences Data Archive \parencite{lsda}. \textcite{meck2009}
describieron ese proyecto y sus signos vitales: al inicio del reposo, el pulso y la presión diastólica bajan de forma
transitoria y hay una diuresis temprana que se estabiliza en unos 3 días. \textcite{platts2009} describieron las
adaptaciones cardiovasculares en 13 sujetos del mismo programa.

\begin{cuidado}
El reposo en cama reproduce el desplazamiento de fluidos y la descarga del cuerpo, pero no la radiación, el CO$_2$ de
la cabina ni el aislamiento de una nave. En vuelo hay efectos que el reposo no muestra, como el aumento de la
temperatura central en misiones largas \parencite{stahn2017} y las paradojas cardiovasculares que revisa
\textcite{norsk2020}. Los datos de vuelo públicos que encontró el equipo, como el estudio de gemelos de la NASA
\parencite{garrettbakelman2019} o el atlas SOMA \parencite{overbey2024}, son análisis en días puntuales, no series
diarias de signos vitales. Por eso este trabajo no valida con vuelo real.
\end{cuidado}

\section{Las campañas}

Se usan tres campañas de la UTMB: la campaña~1, la campaña~3 y el CFT70. En total son \cifra{datos.sujetos} sujetos
(\cifra{datos.c1} de la campaña~1, \cifra{datos.c3} de la campaña~3 y \cifra{datos.cft70} del CFT70) y
\cifra{datos.mediciones} mediciones diarias: \cifra{datos.med.pre} antes del reposo, \cifra{datos.med.cama} en cama y
\cifra{datos.med.post} después. La \cref{tab:campanas} resume cada una.

\begin{table}[htbp]
\caption[Campañas de reposo en cama]{Campañas de reposo en cama. Los días se cuentan desde el primer día en cama
(día~0); el pre-reposo es negativo. Los rangos son los que define cada protocolo.}
\label{tab:campanas}
\begin{tabular}{lrrccc}\toprule
 Campaña & Sujetos & Mediciones & Pre-reposo & En cama & Post-reposo \\
\midrule
Campaña 1 & 3 & 1518 & −10 a −1 & 0 a 59 & 60 a 73 \\
Campaña 3 & 23 & 18\,991 & −16 a −1 & 0 a 91 & 60 a 105 \\
CFT70 & 3 & 1702 & −21 a −1 & 0 a 69 & 70 a 82 \\
\midrule Total & 29 & 22\,211 & & & \\
\bottomrule\end{tabular}

\end{table}

Contando los días con pulso medido, cada sujeto tiene entre \cifra{datos.chk.pulso.pre.min} y
\cifra{datos.chk.pulso.pre.max} días de pre-reposo y entre \cifra{datos.chk.pulso.cama.min} y
\cifra{datos.chk.pulso.cama.max} en cama. Contando cualquier variable, cada uno pasó entre \cifra{datos.chk.cama.min}
y \cifra{datos.chk.cama.max} días en cama. Los fluidos se registran en menos personas: la diuresis en
\cifra{datos.fluidos.diuresis} y el agua ingerida en \cifra{datos.fluidos.agua}.

Dos detalles de diseño de las campañas importan para leer los resultados. En el CFT70, 2 de los 3 sujetos hicieron
ejercicio como contramedida, y la campaña~3 tiene dos grupos. El corrector poblacional del rango esperado los mezcla
con el resto (\cref{cap:discusion}).

\section{Las ocho variables}

La \cref{tab:variables} muestra las ocho variables diarias. Seis están en el contrato de datos de la estación médica de
NOVA (\texttt{nave/estacion/CONTRATO.md}) y usan sus nombres y unidades. La diuresis y el agua ingerida no están en el
contrato; sus rangos válidos son supuestos.

\begin{table}[htbp]
\caption[Variables diarias]{Variables diarias. \enquote{Rango válido}: un valor fuera de ese rango se trata como error
de sensor y se descarta. \enquote{Ruido diario}: mediana entre personas de la desviación de cada una, que escala la
pérdida del pronóstico. \enquote{Escalón}: magnitud de la anomalía sintética que se inyecta para medir la detección.}
\label{tab:variables}
\centering\footnotesize\setlength{\tabcolsep}{4pt}
\begin{tabular}{llllcccrr}\toprule
 Nombre & Variable & Unidad & Rango válido & Contrato & Sujetos & LifeSnaps & Ruido diario & Escalón \\
\midrule
\var{pulso} & Pulso & lpm & 20–250 & sí & 29 & sí & 2,81 & 10 \\
\var{pa\_sis} & Presión sistólica & mmHg & 50–250 & sí & 29 & no & 6,76 & 15 \\
\var{pa\_dia} & Presión diastólica & mmHg & 30–150 & sí & 29 & no & 6,94 & 10 \\
\var{fr} & Frecuencia respiratoria & rpm & 4–60 & sí & 29 & sí & 0,88 & 4,0 \\
\var{temp} & Temperatura & °C & 30–43 & sí & 29 & sí & 0,44 & 0,8 \\
\var{peso} & Peso & kg & 30–200 & sí & 29 & no & 0,60 & 1,5 \\
\var{diuresis} & Diuresis & g/día & 200–8000 & no & 24 & no & 514 & 1000 \\
\var{agua} & Agua ingerida & g/día & 200–10\,000 & no & 23 & no & 262 & 1000 \\
\bottomrule\end{tabular}

\end{table}

El agua ingerida es la más débil de las ocho: sigue una meta diaria del protocolo, así que su \enquote{valor esperado}
es la meta y una desviación es un desvío del protocolo, no fisiología. Se muestra aparte y no entra en los promedios.

\section{LifeSnaps: la Tierra}

LifeSnaps \parencite{yfantidou2022} es un conjunto público de unos cuatro meses de datos de personas que usaron un reloj
Fitbit en su vida diaria, con licencia CC BY 4.0. Después de la limpieza quedan \cifra{datos.ls.personas} personas,
\cifra{datos.ls.diaspersona} días-persona y \cifra{datos.ls.mediciones} mediciones en \cifra{datos.ls.canales}
canales: pulso en reposo, respiración, temperatura de la piel, variabilidad cardíaca, saturación de oxígeno, estrés,
horas de sueño y minutos de actividad.

Los ocho canales entran al pre-entrenamiento del PatchTST, que es univariado y compartido: el modelo aprende la
dinámica diaria de cualquier serie, no de una variable en particular. Con el reposo en cama, LifeSnaps comparte solo
el pulso, la respiración y la temperatura, y medidos con otro instrumento: el reloj en vez de la enfermería, y la piel
en vez del cuerpo. La presión, el peso y los fluidos solo aparecen en Tierra a través del pre-reposo de cada sujeto.
De las \cifra{datos.ls.personas} personas, \cifra{datos.ls.entrenan} entrenan y 7 quedan aparte para la parada
temprana. La idea de usar los relojes para conocer la fisiología personal no es nueva: \textcite{li2017} mostraron
que los sensores que se llevan puestos siguen la fisiología y la actividad de cada persona, y \textcite{mishra2020}
detectaron COVID-19 antes de los síntomas con datos de relojes inteligentes.

\section{Formato común y limpieza}

Las dos fuentes se llevan a un formato largo, con una fila por medición y las columnas fuente, campaña, sujeto,
instrumento, día, fase, variable y valor. En el reposo, el día se cuenta desde el inicio de la inclinación: el
pre-reposo es negativo, el primer día en cama es el~0 y el post-reposo sigue después del último día en cama. Las fases
son \texttt{PRE\_TEST}, \texttt{IN\_TEST} y \texttt{POST\_TEST}. En LifeSnaps, el día se cuenta desde la primera fecha
de cada persona.

Las tablas de la NASA usan nombres de columna distintos en cada versión de sus formularios; el cargador los unifica y
convierte la temperatura de Fahrenheit a Celsius cuando hace falta. De cada celda toma el primer número, y lo que no
es número, como \enquote{Not reported}, queda vacío. Los valores fuera del rango válido se descartan. El código
comprueba que haya a lo sumo una medición por sujeto, variable y día. Para el pronóstico se arman
\cifra{ventanas.ls} ventanas de LifeSnaps, \cifra{ventanas.pre} de pre-reposo y \cifra{ventanas.cama} en cama
(\cref{sec:pronostico}).

\section{Qué cambia en cama}\label{sec:desplazamiento}

Antes de modelar, conviene ver cuánto cambia cada variable. La \cref{fig:desplazamiento} y la \cref{tab:desplazamiento}
comparan, en las mismas personas, la media desde el día 28 en cama con la media del pre-reposo. El cambio se expresa
en unidades y en desviaciones del día a día del pre-reposo, para que las variables se puedan comparar.

\begin{figure}[htbp]
\begin{tikzpicture}
\begin{axis}[informe, width=0.80\linewidth, height=6.4cm, xmin=-3, xmax=1.2,
 ymin=-0.6, ymax=7.6, ytick={0,...,7}, yticklabels={Peso,Agua,Diuresis,PA sistólica,Respiración,Temperatura,PA diastólica,Pulso},
 xmajorgrids, xlabel={cambio medio desde el día 28 en cama, en desviaciones del día a día del pre-reposo},
 axis y line*=left, axis x line*=bottom, y axis line style={draw=none}, ytick style={draw=none},
 clip=false]
\draw[tinta3, densely dashed] (axis cs:0,-0.6) -- (axis cs:0,7.6);
\draw[serieA, line width=1.6pt] (axis cs:-3.0612,0) -- (axis cs:-1.3168,0);
\fill[serieA] (axis cs:-2.1890,0) circle[radius=2.6pt];
\node[anchor=west, font=\footnotesize\sffamily\color{tinta2}] at (axis cs:1.25,0) {−1,31\,kg};
\draw[serieA, line width=1.6pt] (axis cs:-0.7256,1) -- (axis cs:-0.0671,1);
\fill[serieA] (axis cs:-0.3963,1) circle[radius=2.6pt];
\node[anchor=west, font=\footnotesize\sffamily\color{tinta2}] at (axis cs:1.25,1) {−104\,g/día};
\draw[serieA, line width=1.6pt] (axis cs:-0.3950,2) -- (axis cs:0.1639,2);
\fill[serieA] (axis cs:-0.1155,2) circle[radius=2.6pt];
\node[anchor=west, font=\footnotesize\sffamily\color{tinta2}] at (axis cs:1.25,2) {−59\,g/día};
\draw[serieA, line width=1.6pt] (axis cs:-0.3168,3) -- (axis cs:0.2139,3);
\fill[serieA] (axis cs:-0.0514,3) circle[radius=2.6pt];
\node[anchor=west, font=\footnotesize\sffamily\color{tinta2}] at (axis cs:1.25,3) {−0,35\,mmHg};
\draw[serieA, line width=1.6pt] (axis cs:-0.3757,4) -- (axis cs:0.3024,4);
\fill[serieA] (axis cs:-0.0366,4) circle[radius=2.6pt];
\node[anchor=west, font=\footnotesize\sffamily\color{tinta2}] at (axis cs:1.25,4) {−0,04\,rpm};
\draw[serieA, line width=1.6pt] (axis cs:-0.0548,5) -- (axis cs:0.2955,5);
\fill[serieA] (axis cs:0.1203,5) circle[radius=2.6pt];
\node[anchor=west, font=\footnotesize\sffamily\color{tinta2}] at (axis cs:1.25,5) {+0,03\,°C};
\draw[serieA, line width=1.6pt] (axis cs:-0.0750,6) -- (axis cs:0.5044,6);
\fill[serieA] (axis cs:0.2147,6) circle[radius=2.6pt];
\node[anchor=west, font=\footnotesize\sffamily\color{tinta2}] at (axis cs:1.25,6) {+1,49\,mmHg};
\draw[serieA, line width=1.6pt] (axis cs:0.1163,7) -- (axis cs:0.7087,7);
\fill[serieA] (axis cs:0.4125,7) circle[radius=2.6pt];
\node[anchor=west, font=\footnotesize\sffamily\color{tinta2}] at (axis cs:1.25,7) {+2,58\,lpm};
\end{axis}
\end{tikzpicture}

\caption[Cuánto cambia cada variable en cama]{Cambio medio desde el día 28 en cama respecto del pre-reposo, en las
mismas personas, medido en desviaciones del día a día del pre-reposo, con su IC~95\pct{} (t de Student). El peso es
la única variable que se mueve más que su ruido diario.}
\label{fig:desplazamiento}
\end{figure}

\begin{table}[htbp]
\caption[Desplazamiento de cada variable en cama]{Desplazamiento de cada variable: media desde el día 28 en cama menos
media del pre-reposo, pareado por sujeto. \enquote{DE diaria}: mediana entre personas de la desviación del pre-reposo.
\enquote{Post $-$ pre}: media del post-reposo menos la del pre-reposo.}
\label{tab:desplazamiento}
\small
\begin{tabular}{llrrrrrr}\toprule
 Variable & Unidad & Sujetos & Cama $-$ pre & IC 95\,\% & DE diaria & En DE & Post $-$ pre \\
\midrule
Peso & kg & 29 & −1,31 & ±\,0,52 & 0,60 & −2,19 & −0,07 \\
Agua ingerida & g/día & 23 & −104 & ±\,86 & 262 & −0,40 & 15 \\
Diuresis & g/día & 24 & −59 & ±\,144 & 514 & −0,12 & −275 \\
Presión sistólica & mmHg & 29 & −0,35 & ±\,1,79 & 6,76 & −0,05 & −0,19 \\
Frecuencia respiratoria & rpm & 29 & −0,04 & ±\,0,35 & 1,03 & −0,04 & 0,08 \\
Temperatura & °C & 29 & 0,03 & ±\,0,05 & 0,27 & 0,12 & 0,05 \\
Presión diastólica & mmHg & 29 & 1,49 & ±\,2,01 & 6,94 & 0,21 & 0,15 \\
Pulso & lpm & 29 & 2,58 & ±\,1,85 & 6,26 & 0,41 & 7,77 \\
\bottomrule\end{tabular}

\end{table}

El peso baja \cifra{desp.peso}\un{kg} (IC~95\pct{} $\pm$\cifra{desp.peso.ic}), unas \cifra{desp.peso.sd}
desviaciones de su variación de un día a otro, y vuelve casi al nivel previo después del reposo
(\cifra{desp.peso.post}\un{kg}). El pulso sube \cifra{desp.pulso}\un{lpm} ($\pm$\cifra{desp.pulso.ic}) y, al
levantarse, queda \cifra{desp.pulso.post}\un{lpm} sobre el pre-reposo. La presión diastólica sube
\cifra{desp.pa_dia}\un{mmHg}, con un intervalo de $\pm$\cifra{desp.pa_dia.ic} que cruza el cero, y el agua ingerida
baja \cifra{desp.agua}\un{g/día} ($\pm$\cifra{desp.agua.ic}), aunque sigue una meta. Las demás variables cambian
poco: a lo sumo un octavo de su ruido diario.

\begin{clave}
El desplazamiento es grande solo en el peso. En el pulso y la presión, las variables que protagonizarán la corrida
siguiente, el cambio es menor que la mitad de su ruido diario. Eso anticipa el resultado del \cref{cap:resultados}: el
monitor de Tierra falla sobre todo en el peso, y es ahí donde la adaptación gana más. Estas cifras se miden desde el
día 28, así que no ven los cambios del inicio del reposo que describen \textcite{meck2009}.
\end{clave}
